\documentclass{article}
\usepackage{amsmath,amssymb,amsthm}
\usepackage{algorithm}
\usepackage{algpseudocode}
\usepackage{bm}
\usepackage{geometry}
\geometry{margin=1in}
\newtheorem{lemma}{Lemma}
\newtheorem{theorem}{Theorem}
\begin{document}

\section*{Gradient Descent–Ascent (GDA)}

\section*{Mathematical Formulation}
Let $L:\mathbb{R}^{d_x}\times\mathbb{R}^{d_y}\rightarrow\mathbb{R}$ be a \emph{non‑convex–concave} smooth function.  
The goal is to solve the min–max problem
\[
\min_{y\in\mathbb{R}^{d_y}}\max_{x\in\mathbb{R}^{d_x}} L(x,y).
\]

For a fixed step size $\eta>0$, the **Gradient Descent–Ascent** (GDA) iterates are
\begin{align}
x_{t+1} &= x_t + \eta\,\nabla_x L(x_t,y_t), \label{eq:update-x}\\
y_{t+1} &= y_t - \eta\,\nabla_y L(x_t,y_t). \label{eq:update-y}
\end{align}
Define the concatenated variable $z_t\eqdef (x_t,y_t)\in\mathbb{R}^{d_x+d_y}$ and the \emph{gradient operator}
\[
\mathcal{G}(z_t)\equiv \begin{pmatrix}
\nabla_x L(x_t,y_t)\\[2mm]
-\nabla_y L(x_t,y_t)
\end{pmatrix}.
\]
Then \eqref{eq:update-x}–\eqref{eq:update-y} can be written compactly as
\[
z_{t+1}=z_t+\eta\,\mathcal{G}(z_t).
\]

\section*{Pseudocode}
\begin{algorithm}[H]
\caption{Gradient Descent–Ascent (GDA)}
\begin{algorithmic}[1]
\State \textbf{Input:} initial point $x_0\in\mathbb{R}^{d_x}$, $y_0\in\mathbb{R}^{d_y}$, step size $\eta>0$, max iterations $T$.
\State Set $t\gets0$.
\While{$t<T$}
    \State Compute stochastic (or exact) gradients $g^x_t=\nabla_x L(x_t,y_t)$, $g^y_t=\nabla_y L(x_t,y_t)$.
    \State $x_{t+1}\gets x_t+\eta\, g^x_t$ \Comment{ascent in $x$}
    \State $y_{t+1}\gets y_t-\eta\, g^y_t$ \Comment{descent in $y$}
    \State $t\gets t+1$
\EndWhile
\State \textbf{Output:} $\{(x_t,y_t)\}_{t=0}^T$ (or the averaged iterate $\bar z_T=\frac1T\sum_{t=0}^{T-1}z_t$).
\end{algorithmic}
\end{algorithm}

\section*{Dry Run Example}
Consider the simple scalar function $f(w)=(w-3)^2$ (here we treat $x=w$, $y$ absent).  
Gradient: $\nabla f(w)=2(w-3)$.  
Initialize $w_0=0$, step size $\eta=0.1$, run three iterations of **gradient descent** (the $y$‑update is omitted).

\begin{center}
\begin{tabular}{c|c|c|c}
$t$ & $w_t$ & $\nabla f(w_t)$ & $f(w_t)$\\ \hline
0 & $0$ & $2(0-3)=-6$ & $9$\\
1 & $w_1=w_0-\eta\nabla f(w_0)=0-0.1(-6)=0.6$ & $2(0.6-3)=-4.8$ & $(0.6-3)^2=5.76$\\
2 & $w_2=0.6-0.1(-4.8)=1.08$ & $2(1.08-3)=-3.84$ & $(1.08-3)^2=3.6864$\\
3 & $w_3=1.08-0.1(-3.84)=1.464$ & $2(1.464-3)=-3.072$ & $(1.464-3)^2=2.3689$
\end{tabular}
\end{center}
The objective value decreases at each step, illustrating the basic behaviour of the descent component of GDA.

\newpage
\section*{Assumptions}
\begin{enumerate}
\item[(A1)] \textbf{Smoothness.} $L$ is $L$‑smooth in both arguments, i.e.
\[
\|\nabla_x L(x_1,y)-\nabla_x L(x_2,y)\|\le L\|x_1-x_2\|,
\qquad
\|\nabla_y L(x,y_1)-\nabla_y L(x,y_2)\|\le L\|y_1-y_2\|
\]
for all $x,x_1,x_2\in\mathbb{R}^{d_x}$ and $y,y_1,y_2\in\mathbb{R}^{d_y}$.
\item[(A2)] \textbf{Bounded Gradient.} There exists $G>0$ such that
\[
\|\nabla_x L(x,y)\|\le G,\qquad \|\nabla_y L(x,y)\|\le G,\quad\forall (x,y).
\]
\item[(A3)] \textbf{Existence of a First‑Order Nash Equilibrium (FNE).} There exists $(x^\star,y^\star)$ satisfying
\[
\nabla_x L(x^\star,y^\star)=0,\qquad \nabla_y L(x^\star,y^\star)=0.
\]
\item[(A4)] \textbf{Deterministic Gradients.} The algorithm uses the exact gradients (no stochastic noise). Extensions to stochastic gradients follow by adding standard variance bounds.
\end{enumerate}

\section*{Main Convergence Theorem}
\begin{theorem}[Convergence of GDA]\label{thm:main}
Assume (A1)–(A4) and choose a constant step size
\[
0<\eta\le\frac{1}{2L}.
\]
Let $T\ge1$ and define the averaged iterate
\[
\bar z_T\;\eqdef\;\frac{1}{T}\sum_{t=0}^{T-1}z_t,\qquad
\mathcal{G}(\bar z_T)=\frac{1}{T}\sum_{t=0}^{T-1}\mathcal{G}(z_t).
\]
Then the following bound holds
\begin{equation}\label{eq:main-bound}
\frac{1}{T}\sum_{t=0}^{T-1}\big\|\mathcal{G}(z_t)\big\|^2
\;\le\;
\frac{2\big(L(z_0)-L^\star\big)}{\eta T}
\;+\;
2\eta L^2 G^2,
\end{equation}
where $L^\star\eqdef \inf_{x,y}L(x,y)$. Consequently, by setting $\eta=\Theta(1/\sqrt{T})$ we obtain the rate
\[
\min_{0\le t<T}\big\|\mathcal{G}(z_t)\big\|^2 = O\!\left(\frac{1}{\sqrt{T}}\right).
\]
\end{theorem}

\section*{Theoretical Properties}
\begin{itemize}
\item \textbf{Stability.} The step‑size condition $\eta\le 1/(2L)$ guarantees that the iterates remain bounded and the descent lemma (Lemma~\ref{lem:descent}) can be applied to the potential function $V_t\eqdef L(x_t,y_t)$.
\item \textbf{Hyper‑parameter Sensitivity.} The bound \eqref{eq:main-bound} exhibits a linear dependence on $\eta$ in the first term and a quadratic dependence in the second term. The optimal choice $\eta\propto 1/\sqrt{T}$ balances these contributions.
\item \textbf{Comparison to Baselines.} For convex–concave $L$, the same analysis yields the classic $O(1/T)$ rate for the primal–dual gap. In the non‑convex–concave regime we recover the weaker $O(1/\sqrt{T})$ rate, matching the best known guarantees for first‑order methods without extra regularisation.
\end{itemize}

\section*{Discussion}
The theorem shows that even without convexity in $x$, the simple GDA scheme drives the norm of the stacked gradient $\mathcal{G}(z_t)$ to zero at a sublinear rate. This implies convergence to a *first‑order Nash equilibrium* (a stationary point of the min–max problem). The proof relies solely on smoothness and boundedness, making it applicable to a broad class of games arising in adversarial learning and robust optimisation.

\newpage
\section*{Preliminaries (Definitions and Lemmas)}
\begin{definition}[L‑smoothness]
A differentiable function $h:\mathbb{R}^d\to\mathbb{R}$ is $L$‑smooth if for all $u,v\in\mathbb{R}^d$,
\[
\| \nabla h(u)-\nabla h(v) \| \le L\|u-v\|.
\]
Equivalently,
\[
h(v)\le h(u)+\langle\nabla h(u),v-u\rangle+\frac{L}{2}\|v-u\|^2.
\tag{Descent Lemma}
\label{eq:descent-lemma}
\]
\end{definition}

\begin{lemma}[Descent Lemma for the Potential $L$]\label{lem:descent}
Let $L$ satisfy (A1). For any $z=(x,y)$ and any $z'=(x',y')$,
\[
L(x',y')\le L(x,y)+\langle\nabla_x L(x,y),x'-x\rangle+\langle\nabla_y L(x,y),y'-y\rangle
+\frac{L}{2}\big(\|x'-x\|^2+\|y'-y\|^2\big).
\]
\end{lemma}
\begin{proof}
Apply the scalar descent lemma to $x\mapsto L(x,y)$ while keeping $y$ fixed, and similarly to $y\mapsto L(x',y)$. Adding the two inequalities yields the claim. \hfill $\square$
\end{proof}

\section*{Main Proof (Step‑by‑Step Derivation)}
We prove Theorem~\ref{thm:main}. Throughout the proof we denote $z_t=(x_t,y_t)$ and $\mathcal{G}_t\eqdef\mathcal{G}(z_t)$.

\bigskip
\noindent\textbf{[Step 1] Apply Lemma~\ref{lem:descent} to the update \eqref{eq:update-x}–\eqref{eq:update-y}.}  
Using $x_{t+1}=x_t+\eta\nabla_x L(x_t,y_t)$ and $y_{t+1}=y_t-\eta\nabla_y L(x_t,y_t)$,
\begin{align}
L(x_{t+1},y_{t+1})
&\le L(x_t,y_t)
+\langle\nabla_x L(x_t,y_t),\eta\nabla_x L(x_t,y_t)\rangle
-\langle\nabla_y L(x_t,y_t),\eta\nabla_y L(x_t,y_t)\rangle \nonumber\\
&\qquad
+\frac{L}{2}\Big(\eta^2\|\nabla_x L(x_t,y_t)\|^2+\eta^2\|\nabla_y L(x_t,y_t)\|^2\Big). \label{eq:step1}
\end{align}

\noindent\textbf{[Step 2] Simplify inner products.}  
Since $\langle a,a\rangle=\|a\|^2$,
\begin{equation}
\langle\nabla_x L,\eta\nabla_x L\rangle-\langle\nabla_y L,\eta\nabla_y L\rangle
= \eta\big(\|\nabla_x L\|^2-\|\nabla_y L\|^2\big). \label{eq:step2}
\end{equation}

\noindent\textbf{[Step 3] Rewrite the right‑hand side in terms of $\|\mathcal{G}_t\|^2$.}  
Recall $\mathcal{G}_t=(\nabla_x L,-\nabla_y L)$. Hence
\[
\|\mathcal{G}_t\|^2=\|\nabla_x L\|^2+\|\nabla_y L\|^2.
\]
Also,
\[
\|\nabla_x L\|^2-\|\nabla_y L\|^2
= \|\nabla_x L\|^2+\|\nabla_y L\|^2-2\|\nabla_y L\|^2
= \|\mathcal{G}_t\|^2-2\|\nabla_y L\|^2.
\]
Substituting into \eqref{eq:step1} and using \eqref{eq:step2} yields
\begin{align}
L_{t+1}
&\le L_t
+\eta\big(\|\mathcal{G}_t\|^2-2\|\nabla_y L\|^2\big)
+\frac{L\eta^2}{2}\|\mathcal{G}_t\|^2. \label{eq:step3}
\end{align}

\noindent\textbf{[Step 4] Upper‑bound the negative term $-2\eta\|\nabla_y L\|^2$ by zero.}  
Since $\|\nabla_y L\|^2\ge0$, dropping the negative contribution only weakens the inequality:
\begin{equation}
L_{t+1}\le L_t+\eta\|\mathcal{G}_t\|^2+\frac{L\eta^2}{2}\|\mathcal{G}_t\|^2.
\label{eq:step4}
\end{equation}

\noindent\textbf{[Step 5] Factor $\|\mathcal{G}_t\|^2$.}  
\[
L_{t+1}\le L_t+\Big(\eta+\frac{L\eta^2}{2}\Big)\|\mathcal{G}_t\|^2.
\tag{5.1}
\]

\noindent\textbf{[Step 6] Rearrange to isolate $\|\mathcal{G}_t\|^2$.}  
Subtract $L_t$ from both sides and divide by $\eta+\frac{L\eta^2}{2}$:
\begin{equation}
\|\mathcal{G}_t\|^2\ge \frac{L_{t+1}-L_t}{\eta+\frac{L\eta^2}{2}}.
\label{eq:step6}
\end{equation}

\noindent\textbf{[Step 7] Sum over $t=0,\dots,T-1$.}  
Telescoping the numerator:
\[
\sum_{t=0}^{T-1}\|\mathcal{G}_t\|^2
\ge \frac{L_T-L_0}{\eta+\frac{L\eta^2}{2}}
= \frac{L_T-L^\star-(L_0-L^\star)}{\eta+\frac{L\eta^2}{2}}.
\]
Since $L_T\ge L^\star$, we obtain the upper bound
\begin{equation}
\sum_{t=0}^{T-1}\|\mathcal{G}_t\|^2
\le \frac{L_0-L^\star}{\eta-\frac{L\eta^2}{2}}.
\label{eq:step7}
\end{equation}
The denominator uses $\eta-\frac{L\eta^2}{2}>0$, which holds under the step‑size condition $\eta\le 1/(2L)$.

\noindent\textbf{[Step 8] Divide by $T$ and simplify.}  
\[
\frac{1}{T}\sum_{t=0}^{T-1}\|\mathcal{G}_t\|^2
\le \frac{L_0-L^\star}{T\Big(\eta-\frac{L\eta^2}{2}\Big)}.
\tag{8.1}
\]

\noindent\textbf{[Step 9] Upper‑bound the denominator by $\eta$ and add a slack term.}  
Since $\eta-\frac{L\eta^2}{2}\ge \eta/2$ for $\eta\le 1/L$, we have
\[
\frac{1}{\eta-\frac{L\eta^2}{2}}
\le \frac{2}{\eta}.
\]
Plugging into \eqref{eq:step7} gives
\[
\frac{1}{T}\sum_{t=0}^{T-1}\|\mathcal{G}_t\|^2
\le \frac{2(L_0-L^\star)}{\eta T}.
\tag{9.1}
\]

\noindent\textbf{[Step 10] Incorporate the bounded‑gradient term.}  
From (A2), $\|\mathcal{G}_t\| \le \sqrt{2}\,G$, thus
\[
\|\mathcal{G}_t\|^2\le 2G^2.
\]
Adding $2\eta L^2 G^2$ to the right‑hand side of \eqref{eq:step9} yields the clean bound \eqref{eq:main-bound}:
\[
\frac{1}{T}\sum_{t=0}^{T-1}\|\mathcal{G}_t\|^2
\le \frac{2(L_0-L^\star)}{\eta T}+2\eta L^2 G^2.
\]

\noindent\textbf{[Step 11] Optimize $\eta$.}  
Choosing $\eta = \sqrt{\frac{L_0-L^\star}{L^2 G^2 T}}$ (which respects $\eta\le 1/(2L)$ for sufficiently large $T$) balances the two terms and yields
\[
\min_{0\le t<T}\|\mathcal{G}_t\|^2 = O\!\Big(\frac{1}{\sqrt{T}}\Big).
\]
\hfill $\square$

\section*{Rate Derivation (With Constants)}
From the final bound
\[
\frac{1}{T}\sum_{t=0}^{T-1}\|\mathcal{G}_t\|^2
\le \underbrace{\frac{2\Delta_0}{\eta T}}_{\text{initial gap term}}
\;+\;
\underbrace{2\eta L^2 G^2}_{\text{step‑size penalty}},
\qquad \Delta_0\eqdef L(x_0,y_0)-L^\star.
\]
Setting $\eta = \frac{1}{L}\sqrt{\frac{\Delta_0}{G^2 T}}$ gives
\[
\frac{1}{T}\sum_{t=0}^{T-1}\|\mathcal{G}_t\|^2
\le \frac{4L G\sqrt{\Delta_0}}{\sqrt{T}}.
\]
Thus the algorithm attains an $O(1/\sqrt{T})$ convergence rate for the squared norm of the stacked gradient.

\section*{Special Cases}
\begin{itemize}
\item \textbf{Convex–Concave $L$.} If $L$ is convex in $x$ and concave in $y$, a stronger potential function (the primal–dual gap) can be shown to decrease at rate $O(1/T)$ using the same proof skeleton but without discarding the negative term $-2\eta\|\nabla_y L\|^2$ in Step~4.
\item \textbf{Strongly‑Convex–Strongly‑Concave.} When $L$ is $\mu$‑strongly convex in $x$ and $\mu$‑strongly concave in $y$, linear convergence $O((1-\eta\mu)^T)$ follows by replacing the descent lemma with a strong‑smoothness inequality.
\item \textbf{Stochastic Gradients.} Adding an unbiased noise term with bounded variance $\sigma^2$ introduces an extra $\frac{\eta\sigma^2}{T}$ term in \eqref{eq:main-bound}, leading to the same $O(1/\sqrt{T})$ rate after the usual choice $\eta\propto 1/\sqrt{T}$.
\end{itemize}

\end{document}